\chapter{Conclusion}

\section{Summary of Accomplishments}

This project set out to build an end-to-end system that transforms recorded lectures into structured learning materials - transcripts, summaries, and multiple-choice questions - without requiring manual effort from the instructor or the learner. The final implementation achieves that goal across every layer of the stack.

On the infrastructure side, the system runs as a set of containerised services orchestrated with Docker Compose, comprising a Node.js REST API, a React single-page frontend, a FastAPI AI service, MySQL and MongoDB databases, Redis for job queuing, and MinIO for object storage. A BullMQ worker process handles all long-running AI jobs asynchronously, and a CI pipeline on GitHub Actions validates every push. Videos are accepted either by direct upload or by YouTube URL, and the frontend polls automatically until processing completes.

The AI pipeline integrates three models, each chosen for a specific capability. Whisper converts lecture audio to text; we evaluated five model sizes against a 30-minute Yale lecture and selected the \textit{small} variant as the best balance of accuracy (WER 18.68\%, CER 16.47\%) and runtime (114 seconds). For summarization, we adopted a Longformer Encoder-Decoder model fine-tuned on long-document data. Evaluated across three public benchmarks, it achieved a ROUGE-1 of 0.3403 on PubMed biomedical articles — the domain closest to academic lecture text — with an average inference time of 16 seconds per document on CPU. For question generation, a two-stage pipeline produces complete multiple-choice items: a sequence-to-sequence model fine-tuned on reading comprehension data generates the question and answer, and a second model fine-tuned on multiple-choice examination data generates three distractors. A word-embedding fallback covers cases where the distractor model underperforms. Evaluated on 20 SQuAD v1.1 contexts, the pipeline achieved a ROUGE-1 of 0.3528 for question quality, a 90\% generation rate, and 100\% answer-echo avoidance across all generated items.

Together, these components form a working system that a student can use to process a recorded lecture end-to-end in a single session.

\section{Limitations}

Despite meeting its core objectives, the system has several limitations worth acknowledging.

All AI inference runs on CPU inside the Docker container. For summarization, this means 16 seconds per PubMed-length article under optimal conditions; very long lectures or simultaneous users would cause noticeable queuing delays. The system has not been load-tested beyond single-user scenarios.

The summarization model was not fine-tuned on lecture-domain data. Although it performs well on long structured documents, academic speech - with its digressions, hedged assertions, and incomplete sentences — differs from the book and biomedical text the model trained on. The generated summaries are coherent but occasionally miss lecture-specific emphasis.

Distractor quality has room for improvement: 61.1\% of MCQs passed the no-filler check, meaning a substantial minority contained generic or degenerate options. This is most pronounced when the answer span is a short number or proper noun, where the distractor model lacks sufficient topical context.

The system currently processes one question per passage context. Generating multiple non-redundant questions from the same passage - a common requirement in real assessments - is not yet supported.

\section{Future Work}

Several directions are open for continued development.

\textbf{Model acceleration.} Exporting the LED and question-generation models to ONNX and applying dynamic quantization would reduce CPU inference latency substantially. GPU support via an environment variable switch — already prepared in the Dockerfile — would allow deployment on GPU-equipped hosts with no code change.

\textbf{Domain fine-tuning.} Fine-tuning the summarization model on a lecture transcript corpus, or on a mixture of academic speech and long-form documents, would likely improve summary coherence for the target domain. A small labelled dataset of lecture–summary pairs, even a few hundred examples, would be sufficient for supervised continued training.

\textbf{Improved distractor generation.} Filtering the distractor model's output with a named-entity tagger and a subject-specific lexicon would reduce filler options. Alternatively, replacing the word-embedding fallback with a retrieval step over the passage itself — selecting plausible but incorrect spans — would tie distractors more tightly to the source material.

\textbf{Multi-question generation.} Extending the pipeline to extract multiple answer spans per passage and generate a non-redundant question set would allow the system to produce full quiz sheets from a single lecture without manual intervention.

\textbf{Multilingual support.} Whisper already supports over 90 languages, and the translation endpoint is wired into the backend. Extending the summarization and MCQ pipelines to operate on translated text — or replacing them with multilingual sequence-to-sequence models - would open the system to non-English lecture content.

\textbf{Cloud deployment and scalability.} Moving the object store to a managed service, deploying the AI worker on an auto-scaling GPU instance, and placing the API behind a load balancer would allow the platform to serve concurrent users reliably. A CI/CD pipeline already exists; the remaining work is infrastructure provisioning and a production environment configuration.

\textbf{Richer learning features.} Automatic mind-map generation from the transcript, flashcard sets derived from the question–answer pairs, and an AI chatbot that answers questions using the lecture content as a retrieval corpus are all natural extensions that would deepen the system's value as a learning support tool.